\chapter{داده و آرایش آزمایش}
\label{ch:dataset}

\section{پیکره کامل \lr{DocRED}}
\label{sec:corpus}

کار روی این داده دو مرحله داشت. اول سراغ \lr{Re-DocRED} رفتیم، یعنی همان $4{,}053$ سند برچسب‌خورده انسانی \lr{DocRED} با برچسب‌هایی که بازبینی و کامل شده‌اند. بعد که این هم کوچک ارزیابی شد، کل پیکره \lr{DocRED} را برداشتیم؛ چیزی $26$ برابر بزرگ‌تر که بخش عمده‌اش را نه انسان، بلکه هم‌ترازی خودکار با \lr{Wikidata} برچسب زده است. جدول \ref{tab:corpus} ابعاد هر بخش را می‌دهد.

\begin{table}[ht]
\centering
\caption{ابعاد پیکره کامل \lr{DocRED}. اعداد با اسکریپت آمار از روی فایل‌های \lr{Parquet} شمرده شده‌اند و نقل‌قول از مقاله نیستند.}
\label{tab:corpus}
\begin{adjustbox}{max width=\textwidth}
\small
\begin{tabular}{|r|c|c|c|c|c|c|}
\hline
بخش & نوع برچسب & سند & جمله & واژه & موجودیت & سه‌تایی \\
\hline
\lr{train\_annotated} & انسانی & $3{,}053$ & $24{,}256$ & $603{,}468$ & $59{,}493$ & $38{,}180$ \\
\hline
\lr{validation} & انسانی & $998$ & $8{,}057$ & $200{,}077$ & $19{,}528$ & $12{,}275$ \\
\hline
\lr{test} & مخفی & $1{,}000$ & \lr{--} & $197{,}787$ & $19{,}539$ & \lr{--} \\
\hline
\lr{train\_distant} & نظارت‌ازدور & $101{,}873$ & $828{,}115$ & $20{,}328{,}858$ & $1{,}965{,}484$ & $1{,}505{,}638$ \\
\hline
جمع & & $106{,}924$ & & $21{,}330{,}190$ & $2{,}064{,}044$ & $1{,}556{,}093$ \\
\hline
\end{tabular}
\end{adjustbox}
\end{table}

درباره برچسب‌ها دو چیز را باید همین‌جا روشن کرد. برچسب‌های اصلی \lr{DocRED} ناقص‌اند و بخش \ref{sec:two-labels} همین نقص را با یک مجموعه پیش‌بینی ثابت و دو مجموعه برچسب اندازه می‌گیرد. نکته دوم اینکه برچسب بخش نظارت‌ازدور کار انسان نیست، برای همین آنچه روی آن بخش گزارش می‌کنیم «توافق با برچسب ازدور» است و نه $F_1$ به معنای متعارفش.

\section{چرا این داده چالشی است}
\label{sec:why-hard}

\subsection{استدلال چندجمله‌ای اجباری}
$64{,}524$ سه‌تایی، یعنی $53.5\%$ کل رابطه‌ها، میان دو موجودیتی برقرار است که هیچ‌وقت در یک جمله کنار هم نیامده‌اند. ترجمه عملی‌اش این است که یک مدل جمله‌محور، هرقدر هم در سطح جمله خوب کار کند، روی بیش از نیمی از این داده صفر می‌گیرد.

\subsection{نامتوازنی شدید کلاس}
تعداد جفت‌موجودیت نامزد $1{,}584{,}994$ است و از این میان فقط $120{,}664$ جفت واقعاً رابطه دارند، یعنی نرخ مثبت $7.61\%$. به بیان دیگر مدل باید $92.4\%$ موارد را درست رد کند، و این همان مسئله مهار توهم است.

\subsection{فضای برچسب بزرگ و دنباله‌دار}
$96$ رابطه وجود دارد، ولی پنج رابطه نخست $52.6\%$ نمونه‌ها را می‌گیرند و هفت رابطه کمتر از $100$ نمونه دارند.

\subsection{چندبرچسبی بودن}
$24{,}636$ جفت‌موجودیت هم‌زمان بیش از یک رابطه دارند، پس ارزیابی ناچار باید چندبرچسبی باشد.

\subsection{هم‌مرجعی و حل موجودیت}
$105{,}688$ ذکر به $78{,}822$ موجودیت نگاشت می‌شود و $18.6\%$ موجودیت‌ها بیش از یک ذکر دارند.

\section{آرایش آزمایش}
\label{sec:setup}

همان آرایش استانداردی را نگه داشتیم که مقاله‌های این بنچمارک استفاده می‌کنند: فهرست موجودیت‌های هر سند به مدل داده می‌شود و مدل باید رابطه‌های بین آن‌ها را پیدا کند و برای هرکدام جمله شاهد را نقل کند. هیچ بخشی از خط لوله فاز اول برای این کار بازنویسی نشد و تنها نقطه اتصال تازه، خواندن هستان‌شناسی از فایل بود.

\section{اعتبارسنجی خود ارزیاب}
\label{sec:evaluator}

قبل از اینکه به هیچ عددی اعتماد کنیم، اول خود ارزیاب را آزمایش کردیم. شش تست آفلاین داریم که نه به شبکه نیاز دارند و نه به مدل زبانی: تست اوراکل، تست نصف بازخوانی، تست موجودیت ناشناخته، تست شکل ظاهری متفاوت، تست مجموعه $\mathrm{Ign}$ و تست سه قاعده بستار.

یک نکته تست اوراکل را باید جدی گرفت. پاسخ کاملاً درست، $F_1$ برابر $99.62$ روی زیرمجموعه $50$ سند و $99.76$ روی کل $500$ سند گرفت، نه $100$. علتش این است که در چند سند دو موجودیت متفاوت نام و نوع یکسان دارند و با نام از هم جدا نمی‌شوند. پس سقف هر سامانه نام‌محور روی این داده حدود $99.7$ است و هر عددی که پایین‌تر می‌آید باید کنار همین سقف خوانده شود.

\section{سامانه دقیقاً از چه نظارتی استفاده می‌کند}
\label{sec:supervision}

این بخش را نوشتیم تا ادعای بزرگ‌تر از واقعیت نکنیم. سامانه هیچ مرحله آموزش گرادیانی ندارد و هیچ وزنی در آن به‌روز نمی‌شود، اما «بدون نظارت» خواندنش هم درست نیست، چون دو چیز از بخش \lr{Train} گرفته‌ایم: هستان‌شناسی با $654$ ترکیب مجاز، و دو سند نمونه با پاسخ کامل طلایی که داخل پرامپت به‌عنوان مثال کاری می‌نشینند.

این کار یک پیامد ناخوشایند دارد. چون هستان‌شناسی هر ترکیبی را که در \lr{Train} دیده نشده ممنوع می‌کند، اگر ترکیبی فقط در \lr{Dev} آمده باشد سامانه هیچ راهی برای پیش‌بینی‌اش ندارد. برای اینکه بدانیم این سقف چقدر گران تمام می‌شود اندازه‌اش گرفتیم: از $17{,}284$ حقیقت طلایی بخش \lr{Dev} فقط $37$ مورد، یعنی $0.21\%$، ترکیب نوعی دارند که در \lr{Train} نیامده است. پس سقف بازخوانی تحمیل‌شده $99.79\%$ است.

معامله‌ای که این فیلتر پیشنهاد می‌کند روشن است: $20.3\%$ خروجی مدل زبانی را می‌گیرد و در عوض $0.21\%$ بازخوانی را می‌سوزاند. با این نسبت، نگه داشتنش تصمیم سختی نبود.

\section{تنظیم روی \lr{Dev} و حساسیت تطبیق}
\label{sec:tuning}

پرامپت‌ها در شش تکرار شکل گرفتند و هر تکرار روی همان بخش \lr{Dev} سنجیده شد که عدد نهایی هم از آن گزارش می‌شود. تصمیم‌هایی مثل اینکه اجماع بین کدام گذرها باشد یا کدام سه قاعده بماند، بعد از دیدن نتیجه \lr{Dev} گرفته شد و نه قبل از آن. برچسب بخش \lr{Test} عمومی نیست، پس در مرحله سوم همین پیکربندی را روی $3{,}053$ سند بخش \lr{Train} برچسب‌خورده اجرا کردیم که در تنظیم پرامپت نقشی نداشت؛ نتیجه در بخش \ref{sec:other-splits} آمده است.

نگرانی دوم، آسان‌گیری تطبیق نام بود. ارزیاب پاسخ را با نام به موجودیت طلایی نگاشت می‌کند و در حالت پیش‌فرض حتی وقتی نوع پیش‌بینی‌شده با نوع طلایی نمی‌خواند باز هم تطبیق را می‌پذیرد. برای اینکه این نگرانی در حد حدس نماند، یک حالت سخت‌گیرانه اضافه کردیم که هر پیش‌بینی با نوع نادرست را دور می‌ریزد. اختلاف دو حالت $0.02$ واحد شد: $28.64$ در برابر $28.66$. پس آسان‌گیری تطبیق عدد را متورم نکرده بود.
